\section{Scripture, the Fathers, and the Formation of Ecclesial Order}

For Christian theology, Scripture establishes creation by a personal God, covenantal command, sin and grace, fulfillment of the Law in Christ, the New Law of the Spirit, ecclesial authority, and a supernatural end beyond temporal peace. These claims transform the questions inherited from Greek philosophy and Roman law. The result is neither a biblical code for every polity nor a merely philosophical account with Christian terminology.

\subsection{Creation, moral knowledge, and revealed law}

Romans 2:14--16 describes Gentiles who do not possess the Mosaic Law yet do by nature what the Law requires, with conscience bearing witness (\sourceurl{https://bible.usccb.org/bible/romans/2}{Romans 2}). The passage has become a principal text in Catholic natural-law reflection. Its precise rhetorical subjects remain disputed. For this study's limited purpose, it supports the claim that moral accountability is not confined to those who possess Israel's written Torah.

Creation, the image of God, wisdom literature, and prophetic appeals to justice also support an intelligible moral order. Biblical and philosophical uses of nature nevertheless require their own exegesis. Creation remains good and providential; human reason and desire are wounded by sin; and human persons are ordered to communion with God.

Covenantal commands furnish a central case of norms known through divine historical revelation. The Decalogue also states central natural-law precepts. A prohibition can therefore be naturally grounded and positively revealed without the two modes becoming identical. Revelation confirms, specifies, and places the precept within covenant and worship.

Jeremiah promises a covenant in which the law is written upon the heart (\sourceurl{https://bible.usccb.org/bible/jeremiah/31}{Jeremiah 31:31--34}). Jesus speaks of fulfilling the Law (\sourceurl{https://bible.usccb.org/bible/matthew/5}{Matthew 5:17--48}), concentrates it in love of God and neighbor (\sourceurl{https://bible.usccb.org/bible/matthew/22}{Matthew 22:34--40}), and gives a new commandment shaped by his self-gift (\sourceurl{https://bible.usccb.org/bible/john/13}{John 13:34--35}). In Aquinas's synthesis, the New Law is principally the grace of the Holy Spirit and secondarily written teaching. The earlier covenant's moral, ceremonial, and judicial precepts therefore require a salvation-historical account of fulfillment; their former juridical form did not cease through human repeal of divine law.

\subsection{Acts 15: a first obligation account}

Acts 15 places a concrete person before several claimed norms. Gentile converts enter a community formed by Israel's Scriptures and Christ's saving work. Some believers argue that circumcision and observance of the Mosaic law are necessary. Apostles and elders deliberate, hear testimony, interpret God's action and Scripture, distinguish what should be imposed, and send a written determination to affected communities (\sourceurl{https://bible.usccb.org/bible/acts/15}{Acts 15}).

The decision combines theological judgment, apostolic competence, a moral minimum, disciplinary provisions directed to communion, and pastoral reception. The Pontifical Biblical Commission reads the event as distinguishing enduring moral requirements from temporary compromise while relating communal discernment to personal conscience (\sourceurl{https://www.vatican.va/roman_curia/congregations/cfaith/pcb_documents/rc_con_cfaith_doc_20080511_bibbia-e-morale_en.html}{\emph{Bible and Morality} 151--153}). The passage does not supply the later law taxonomy, but it displays its practical problem: the source and scope of a claimed divine norm, competent ecclesial judgment, the status of the persons addressed, community unity, and the action now required.

\subsection{Civil authority and its moral limit}

Romans 13:1--7 and 1 Peter 2:13--17 direct Christians to obey governing authority in its service of order and good (\sourceurl{https://bible.usccb.org/bible/romans/13}{Romans 13}; \sourceurl{https://bible.usccb.org/bible/1peter/2}{1 Peter 2}). Acts 5:29 places that obedience under the prior claim of God when authority forbids apostolic witness (\sourceurl{https://bible.usccb.org/bible/acts/5}{Acts 5}). Read together, the texts affirm real political authority and a moral limit to its commands.

Jesus' instruction concerning Caesar and God (Matt 22:15--22) distinguishes claims without assigning every mixed matter to a forum. Early martyr accounts and apologies show Christians praying for rulers and refusing acts of idolatrous worship. Tertullian's \emph{Apology} 33--34, for example, honors the emperor as emperor while denying him divine status (\sourceurl{https://www.newadvent.org/fathers/0301.htm}{text}). The duty of civil loyalty and the duty of religious refusal fell upon the same person.

\subsection{Ecclesial norms before imperial recognition}

The New Testament depicts binding and loosing, apostolic judgment, appointment of ministers, community discipline, adjudication, and conciliar decision (Matt 16 and 18; Acts 15; 1 Cor 5--6; the Pastoral Epistles). These sources ground the Church's capacity to establish binding internal norms and adjudicate disputes before imperial recognition. The 1983 CIC's promulgating constitution later places canonical discipline within that ecclesiological history (\sourceurl{https://www.vatican.va/content/john-paul-ii/en/apost_constitutions/documents/hf_jp-ii_apc_25011983_sacrae-disciplinae-leges.html}{\emph{Sacrae disciplinae leges}}).

The \emph{Didache} joins catechesis, the two ways of life and death, worship, appointment, and community discipline (\sourceurl{https://www.newadvent.org/fathers/0714.htm}{text}). During the first centuries, local and regional councils, episcopal letters, liturgical practice, and penitential discipline generated bodies of ecclesial norms with different territorial scope and later reception. An exhortation, a local canon, and a generally received rule therefore require separate identification.

\subsection{Augustine: eternal measure, temporal peace, and coercion}

In \emph{De libero arbitrio} I.5--6, Augustine makes eternal law the measure of temporal justice and explains why temporal law does not punish every disordered desire. \emph{City of God} XIX.17 presents the pilgrim city as obeying diverse laws that maintain earthly peace insofar as they do not impede worship of God (\sourceurl{https://www.newadvent.org/fathers/120119.htm}{Book XIX.17}). His two cities are communities defined by their loves, not institutional synonyms for Church and state.

Augustine's judgment about coercion changed during the Donatist controversy. Letter 93.16--17 says that he had opposed coercion into unity but revised his position after observing imperial pressure. Letter 185.23--26 defends fines and exile as corrective while rejecting capital punishment in that setting and acknowledging that goodness cannot be produced without the will. The development records a genuine tension among truth, baptismal obligation, public order, persuasion, and coercive means. It later became part of the history against which Catholic teaching on religious freedom developed.

\subsection{Christian empire and diverging canonical trajectories}

Imperial toleration and patronage altered the institutional setting. Councils issued canons; bishops governed communities and sometimes adjudicated civil disputes; emperors convoked councils, enacted laws affecting ecclesiastical institutions, enforced settlements, and sometimes attempted to control doctrine. Marriage, property, clerical status, asylum, worship, and public order could acquire both ecclesiastical and civil consequences. The same persons and acts thus entered distinct but interacting forums.

In the East, conciliar canons, imperial legislation, and later nomocanonical collections placed ecclesiastical and civil texts in coordinated corpora. Justinian's Novel 6 (535) distinguishes priesthood, which serves divine matters, from imperial authority, which administers human affairs, while the enactment itself regulates episcopal qualifications and Church discipline (\sourceurl{https://droitromain.univ-grenoble-alpes.fr/Anglica/N6_Scott.htm}{older English witness}). The text illustrates cooperation and overlap rather than a modern allocation of wholly separate systems.

The Council in Trullo (691--692) received a defined corpus of earlier canons and added extensive discipline; its reception differed between East and West. Later Eastern commentators and collections developed within those sources. The promulgating constitution of the CCEO expressly locates the current Eastern code within this first-millennium conciliar inheritance and the distinct patrimonies of the Eastern Churches (\sourceurl{https://www.vatican.va/content/john-paul-ii/la/apost_constitutions/documents/hf_jp-ii_apc_19901018_sacri-canones.html}{\emph{Sacri canones}}). The CCEO is therefore the result of a different canonical history, not an Eastern translation of Latin law.

In the Latin West, conciliar canons, papal decretals, patristic authorities, penitential books, Roman legal materials, and local collections circulated without one comprehensive code. Reform disputes, the recovery of Justinianic texts, and the growth of schools and papal adjudication eventually created the setting in which Gratian and the decretists worked. The history proceeds from a plurality of sources and forums to juristic efforts at concord, not from an original complete taxonomy to later application.
